\section{Machine-learning searches for the Lax pairs}
\label{sec:ml-search}

We determine coefficients in finite Lax ans\"atze by minimizing an
off-shell residual, following the optimization approach of
Ref.~\cite{KrippendorfLustSyvaeri2021}.  The input data are local values
of the spin and its derivatives,
$(\boldsymbol S,\partial_x\boldsymbol S,
\partial_x^2\boldsymbol S,\partial_t\boldsymbol S)$.
These are sampled subject to the unit-spin constraint and its differentiated
identities, with the time derivative independent of the equations of motion.
The sampling distributions are specified in Section~\ref{sec:ml-numerics}.

For each model $m\in\{5,6\}$, a run holds the deformation parameter
$\alpha_m$ and spectral parameter $\lambda$ fixed.  The equations of motion
and the tensor structures allowed in the Lax ansatz are also supplied.
The trainable parameters are the coefficients of those structures: thirteen
real coefficients for Class 5 and eight for Class 6, each shared by all
samples in a run.  They are adjusted so that the curvature agrees with the
equation-of-motion residual multiplied by the coefficient map determined
from the spatial ansatz.  This yields a candidate Lax pair at the chosen
$(\alpha_m,\lambda)$.  Training requires no target values of a Lax pair.
The numerical runs use real $\alpha_m$ and $\lambda\in\ii\mathbb R_{>0}$,
for which the coefficient relations below are real.  Their analytic verification applies
to complex $\lambda\ne0$.

To write the ans\"atze, represent traceless matrices by three-component
vectors.  With the Pauli matrices acting on the two-dimensional auxiliary
space, define
\begin{equation}
\iota(\boldsymbol v):=-\frac{\ii}{2}v^i\sigma^i,
\qquad
[\iota(\boldsymbol v),\iota(\boldsymbol w)]
=\iota(\boldsymbol v\times\boldsymbol w).
\label{eq:ml-iota}
\end{equation}
Write the candidate matrices as
$U_m^{\rm fit}=\mu_m\Id_2+\iota(\boldsymbol u_m)$ and
$V_m^{\rm fit}=\nu_m\Id_2+\iota(\boldsymbol v_m)$.
Here $\boldsymbol u_m,\boldsymbol v_m$ contain the trainable coefficients,
while the scalar parts $\mu_m,\nu_m$ are prescribed functions, specified
below.  The superscript distinguishes the fitted candidate from the
analytic matrices in the main text.
The vector part of its curvature is
\begin{equation}
\boldsymbol{\mathcal F}^{(m)}
:=\partial_t\boldsymbol u_m-\partial_x\boldsymbol v_m
  +\boldsymbol u_m\times\boldsymbol v_m.
\label{eq:ml-vector-curvature}
\end{equation}
The vectors $\boldsymbol E^{(5)}$ and $\boldsymbol E^{(6)}$ denote
$\partial_t\boldsymbol S$ minus the right-hand side of
Eqs.~\eqref{eq:class5-vector-eom} and \eqref{eq:class6-vector-eom},
respectively.  They are computed from the input data at each sampled point.
All spin contractions and cross products are complex bilinear.
In the losses, $\|\boldsymbol v\|^2=\sum_i|v^i|^2$ is the positive
Euclidean norm squared, and angle brackets denote the arithmetic mean
over the training samples.

\subsection{Class 5}
\label{sec:ml-class5-paper}

The constant matrix $M_5$ in Eq.~\eqref{eq:class5-M-definition} can be
written as $M_5\boldsymbol v=\boldsymbol m_5\times\boldsymbol v$, where
\begin{equation}
\boldsymbol m_5=(\ii,-1,0)^{\mathsf T},
\qquad \boldsymbol m_5^{\mathsf T}\boldsymbol m_5=0,
\qquad M_5^2=\boldsymbol m_5\boldsymbol m_5^{\mathsf T}.
\label{eq:ml-null-vector}
\end{equation}
The ansatz uses the spatial differential operator
$\mathcal D_x\boldsymbol v:=\partial_x\boldsymbol v
+(\alpha_5/2)M_5\boldsymbol v$ and the constant complex orthogonal matrix
\begin{equation}
O_5(\lambda):=\exp(\alpha_5\lambda M_5)
=\Id_3+\alpha_5\lambda M_5
 +\frac{\alpha_5^2\lambda^2}{2}M_5^2,
\qquad O_5^{\mathsf T}O_5=\Id_3,
\qquad\det O_5=1.
\label{eq:ml-O5}
\end{equation}
The spatial ansatz is
\begin{equation}
\boldsymbol u_5
=a_5O_5\boldsymbol S
+b_5\boldsymbol m_5(\boldsymbol m_5^{\mathsf T}\boldsymbol S)
+c_5\boldsymbol m_5,
\label{eq:ml-U5-ansatz}
\end{equation}
and the time ansatz contains the twelve structures
\begin{align}
\boldsymbol v_5={}&
 d_1 O_5(\boldsymbol S\times\mathcal D_x\boldsymbol S)
+d_2\boldsymbol m_5\bigl[\boldsymbol m_5^{\mathsf T}
 (\boldsymbol S\times\mathcal D_x\boldsymbol S)\bigr]
+d_3O_5\boldsymbol S
+d_4\boldsymbol m_5(\boldsymbol m_5^{\mathsf T}\boldsymbol S)
+d_5\boldsymbol m_5
\nonumber\\
&+d_6O_5\mathcal D_x\boldsymbol S
+d_7\boldsymbol m_5(\boldsymbol m_5^{\mathsf T}\mathcal D_x\boldsymbol S)
+d_8\boldsymbol m_5(\boldsymbol m_5^{\mathsf T}\boldsymbol S)^2
+d_9\boldsymbol m_5(\boldsymbol m_5^{\mathsf T}\boldsymbol S)
  (\boldsymbol m_5^{\mathsf T}\mathcal D_x\boldsymbol S)
\nonumber\\
&+d_{10}O_5(\boldsymbol S\times\boldsymbol m_5)
+d_{11}\boldsymbol m_5\bigl[\boldsymbol m_5^{\mathsf T}
(\boldsymbol S\times\boldsymbol m_5)\bigr]
+d_{12}O_5\boldsymbol m_5.
\label{eq:ml-V5-ansatz}
\end{align}
The fitted coefficients are
$(a_5,b_5,c_5,d_1,\ldots,d_{12})$.  The scalar parts are fixed to
\begin{align}
\mu_5&=\frac{1}{4\lambda}+\frac{\alpha_5}{4}S^+,
\nonumber\\
\nu_5&=\frac{\ii\alpha_5}{2}
(S^+\partial_xS^3-S^3\partial_xS^+)
+\frac{\ii\alpha_5^2}{4}(S^+)^2,
\label{eq:ml-class5-scalars}
\end{align}
and obey
$\partial_t\mu_5-\partial_x\nu_5=(\alpha_5/4)E^{(5)}_+$,
where $E^{(5)}_+$ is defined in Eq.~\eqref{eq:class5-residual-plus}.
This identity holds independently of the fitted coefficients, so the
numerical search determines the traceless parts of the pair.

The off-shell loss is
\begin{equation}
\mathcal L_5
=\left\langle
\left\|
\boldsymbol{\mathcal F}^{(5)}
-Q_5\boldsymbol E^{(5)}
\right\|^2
\right\rangle
+\gamma_5\sum_{r\in\mathcal R_5}|r|,
\label{eq:ml-class5-loss}
\end{equation}
where $Q_5$ is the coefficient of $\partial_t\boldsymbol S$ induced by
Eq.~\eqref{eq:ml-U5-ansatz}, $\mathcal R_5$ contains the eight coefficients
$d_5,\ldots,d_{12}$, and $\gamma_5=10^{-7}$.  This is the objective used in
the nine stored Class 5 runs.

The recovered sparse branch has
\begin{equation}
 c_5=d_5=\cdots=d_{12}=0,
\label{eq:ml-class5-sparse}
\end{equation}
while the remaining coefficients reproduce the analytic pair in
Section~\ref{sec:class5-model}.  The numerical values and run-to-run
statistics are given in the repository output files described in
Section~\ref{sec:ml-numerics}.

\subsection{Class 6}
\label{sec:ml-class6-paper}

For Class 6, let $N_6$ be the nilpotent matrix in
Eq.~\eqref{eq:class6-N-definition}.  The search uses the exact finite basis
$\{\Id_3,N_6,N_6^2\}$.  The spatial and time ans\"atze are
\begin{align}
\boldsymbol u_6&=(a_0\Id_3+a_1N_6+a_2N_6^2)\boldsymbol S,
\label{eq:ml-U6-ansatz}\\
\boldsymbol v_6&=(b_0\Id_3+b_1N_6+b_2N_6^2)
(\boldsymbol S\times\partial_x\boldsymbol S)
+(c_0\Id_3+c_1N_6)\boldsymbol S.
\label{eq:ml-V6-ansatz}
\end{align}
The eight trainable coefficients are
$(a_0,a_1,a_2,b_0,b_1,b_2,c_0,c_1)$.  The scalar parts are fixed to
\begin{equation}
\mu_6=\frac{1}{4\lambda},
\qquad
\nu_6=0.
\label{eq:ml-class6-scalars}
\end{equation}

The Class 6 objective can be written off shell as
\begin{equation}
\mathcal L_6
=\left\langle
\left\|
\boldsymbol{\mathcal F}^{(6)}-A_6\boldsymbol E^{(6)}
\right\|^2
\right\rangle,
\qquad
A_6:=a_0\Id_3+a_1N_6+a_2N_6^2.
\label{eq:ml-class6-loss}
\end{equation}
For any fitted coefficients, the time-derivative terms cancel in
$\boldsymbol{\mathcal F}^{(6)}-A_6\boldsymbol E^{(6)}$.
The original stored runs used the on-shell form obtained by substituting the
Class 6 equations of motion for $\partial_t\boldsymbol S$; the two losses are
exactly equal for the same spatial samples and coefficients.  The active
script uses Eq.~\eqref{eq:ml-class6-loss} and also reruns the original
on-shell objective as a convention check.

The recovered coefficients are equivalent to the matrix-square-root form
\begin{equation}
A_6(\lambda)
=\frac{1}{2\lambda}
\left(\Id_3+2\alpha_6\lambda^2N_6
+2\alpha_6^2\lambda^4N_6^2\right),
\label{eq:ml-A6}
\end{equation}
with the corresponding time coefficients fixed by the zero-curvature
condition.  Substitution gives the analytic Class 6 matrices in
Section~\ref{sec:class6-model}.

\subsection{Numerical settings and verification}
\label{sec:ml-numerics}

For Class 5, each run uses $100$ sampled configurations and optimizes with
Adam for $5000$ steps at learning rate $10^{-2}$, followed by L-BFGS for at
most $500$ iterations.  Nine independent seeds are stored.  For Class 6,
each run uses $128$ samples, Adam for $4000$ steps at learning rate
$5\times10^{-3}$, followed by L-BFGS for at most $500$ iterations.  Six
independent seeds are stored.  The samples are generated by drawing Gaussian
vectors, normalizing the spin to unit length, projecting the first derivative
to the tangent plane, and projecting the second derivative so that the second
derivative of the unit-spin constraint is satisfied.  The time derivative is
sampled independently and projected to the tangent plane for the off-shell
losses.

The stored Class 5 runs reproduce the nonzero analytic coefficients with
relative errors below $10^{-5}$ after fixing the common normalization, while
the eight additional candidate coefficients are below $10^{-6}$ in absolute
value.  The stored Class 6 runs reproduce the matrix-square-root relations
within the same tolerance.  The active verification script checks the loss
and its gradient at arbitrary coefficients, verifies the equivalence of the
two Class 6 objectives, reruns the stored seeds, and compares the recovered
matrices with the analytic expressions used in the manuscript.
